% Round 15. New incidence argument. Conditional on the separately proved
% arbitrary-base P1 PWS integration and its algebraic base-change property.
% No history file is modified by this section.
\section{The actual Aeppli formula for projective bundles of every rank}
\label{r15i:section}

The relative diagonal of a higher-dimensional projective space is not
a divisor.  The construction below replaces it by the divisor given
by a marked point on a varying projective line.  Its integration maps
are a tower of projective-bundle projections, for which ordinary PWS
outputs have been established separately.  The birational incidence
map is used only to compare integrals of already defined top forms.

\begin{definition}[The natural projective-line input]
\label{r15i:HPnat}
Write $(\mathrm{HP}_{\mathrm{nat}})_K$ for ordinary PWS integration
along $U\times\mathbf P^1\to U$ in every bidegree, with all ordinary
smooth compact-test pairings, as in Theorem~\ref{r15hp:projection},
together with its naturality for algebraic morphisms between smooth
integral bases, as in Proposition~\ref{r15hp:base-change}.
The assertion is required on ordinary base opens; a projective base
is not required for that analytic input.  Its extension to projective
bundles in Theorem~\ref{r15ps:projective-bundle}, with algebraic
base change \eqref{r15ps:bundle-base-change} and the projection
formula against smooth target forms, follows from this input and
the frozen native finite degree formula.  Only that smooth target
projection formula is needed for the main proof below.
\end{definition}

The name distinguishes this naturality hypothesis from the smaller
product-base-change clause sufficient for the rank-two proof.
No fixed-kernel output hypothesis beyond the PWS integration theorem
is needed in this incidence argument.  We use smooth metrics on line
bundles, whose existence and comparison with model metrics were
proved in Lemma~\ref{r15b:smooth-line-metrics}.  No smooth norm on a
vector bundle of rank greater than one is assumed.

\begin{lemma}[Relative top Chern normalization]
\label{r15i:normalization}
Assume $(\mathrm{HP}_{\mathrm{nat}})_K$.  Let
$a:P=\mathbf P_B(V)\to B$ be a projective bundle of relative
dimension $s$, where $B/K$ is smooth integral and projective.
Let $\zeta$ be the curvature of a smooth metric on
$\mathcal O_P(1)$.  Then the ordinary PWS output satisfies
\[
                         a_*(\zeta^s)=1.
\tag{I15.1}\label{r15i:relative-mass}
\]
For $s=0$ the assertion has the evident identity interpretation.
\end{lemma}

\begin{proof}
The frozen projective-bundle degree lemma is
\texttt{r10c:projective-degree}, formula (R10.1), in
\texttt{round12/inherited/r10\_projective\_degree.tex}.
It states, for any PWS metric on $\mathcal O_P(1)$ and every
compact actual top test $\eta$ on the base, that
\[
                 \int_P\zeta^s a^*\eta=\int_B\eta.
\]
Its proof uses successive coordinate hyperplane sections and
Mihatsch Corollary 5.6 against separately closed actual tests;
it assumes no form-valued projective-bundle pushforward.
The new PWS output has exactly this pairing for every smooth
top test.  Current injectivity therefore gives
\eqref{r15i:relative-mass}.  This also proves the normalization
over a closed residue field $F$ viewed as a $K$-variety: both
sides have native mass $[F:K]$.  No scalar extension
normalization or comparison of different test spaces occurs.
\end{proof}

Fix now a smooth integral projective $K$-variety $S$, a vector
bundle $N$ of rank $r+1$, with $r\geq1$, and the convention of
lines in
\[
                     \pi:E=\mathbf P_S(N)\longrightarrow S.
\]
Put $m=\dim E$.  Denote by $L_y=\mathcal O_E(-1)$ the
tautological line and by $Q=\pi^*N/L_y$ the quotient bundle of
rank $r$.  Define
\[
 \lambda:F=\mathbf P_E(Q)\longrightarrow E,
 \qquad T_F=\mathcal O_F(-1)\subset\lambda^*Q.
\tag{I15.2}\label{r15i:flag-base}
\]
Let $M\subset\lambda^*\pi^*N$ be the inverse image of $T_F$;
it is the rank-two vector bundle in
\[
                0\longrightarrow\lambda^*L_y
                  \longrightarrow M\longrightarrow T_F
                  \longrightarrow0.
\tag{I15.3}\label{r15i:rank-two-subspace}
\]
Finally put
\[
 \rho:I=\mathbf P_F(M)\longrightarrow F,
 \qquad p=\lambda\rho:I\longrightarrow E.
\tag{I15.4}\label{r15i:incidence-tower}
\]
The inclusion $M\subset\lambda^*\pi^*N$ defines a second
map $z:I\to E$.  A point of $F$ specifies a point $y$ of a
fiber of $E\to S$ and a projective line through $y$; a point
of $I$ additionally specifies a point $z$ on that line.
All four spaces $S,E,F,I$ are smooth integral projective over
the original $K$.

\begin{lemma}[The incidence divisor and its Green scalar]
\label{r15i:divisor}
The tautological line map
$\mathcal O_I(-1)\to\rho^*T_F$ cuts out a smooth section
$\Sigma\simeq F$ of $\rho$.  Its divisor line and its normal
line are, respectively,
\[
 \mathcal O_I(\Sigma)
        =z^*\mathcal O_E(1)\otimes\rho^*\mathcal O_F(-1),
 \qquad
 N_{\Sigma/I}
        =\lambda^*\mathcal O_E(1)\otimes\mathcal O_F(-1).
\tag{I15.5}\label{r15i:section-lines}
\]
On $\Sigma$, both $z$ and $p$ are $\lambda$.
The map
\[
        f=(z,p):I\longrightarrow E\times_S E
\tag{I15.6}\label{r15i:incidence-map}
\]
is an isomorphism away from the relative diagonal, whose inverse
image is $\Sigma$.

Choose smooth line metrics with curvatures
$\xi=c_1(\mathcal O_E(1))$ and
$\nu=c_1(\mathcal O_F(1))$.  If $s_\Sigma$ is the displayed
tautological section, set $g=-\log\|s_\Sigma\|$ for the
tensor metric in \eqref{r15i:section-lines}.  With pullbacks
of $\nu$ suppressed, it satisfies
\[
                 Dg=z^*\xi-\nu-\delta_\Sigma
\tag{I15.7}\label{r15i:scalar-Green}
\]
against every compact actual test of the required degree.
The scalar $g$ and its first derivatives are ordinary PWS off
$\Sigma$.
\end{lemma}

\begin{proof}
The kernel of $M\to T_F$ is the line $\lambda^*L_y$.
Consequently the zero locus of the induced map from the
tautological line of $\mathbf P_F(M)$ is exactly its section
$\mathbf P_F(\lambda^*L_y)$, with multiplicity one.
Since $\mathcal O_I(1)=z^*\mathcal O_E(1)$, its divisor
line is the first line in \eqref{r15i:section-lines}.
Restricting to the section gives the second; it is also
$\operatorname{Hom}(\lambda^*L_y,T_F)$, the usual relative
tangent line at that section.  This fixes in particular the
minus sign on $\nu$.

For distinct lines $L_z,L_y\subset N_s$, their span is the
unique two-dimensional subspace containing both.  The image of
$L_z$ in $N_s/L_y$ is nonzero and gives the corresponding line
of $Q_y$.  This construction is algebraic on the complement
of the diagonal and is inverse to \eqref{r15i:incidence-map}.
The section is precisely its inverse image of the diagonal.
In particular the map has native generic degree one.  The
argument uses only the isomorphism on this open; no general
actual pushforward along $f$ is introduced.

The metric is smooth on its line bundle.  Its curvature is
$z^*\xi-\rho^*\nu$.  The ordinary actual Poincare--Lelong
formula, Mihatsch Corollary 5.6, proves
\eqref{r15i:scalar-Green} with all actual compact tests.
Locally off the section, $g$ is a smooth metric weight plus
the valuation of an analytic unit, so it and its first
derivatives are PWS there.
\end{proof}

Define the smooth separately closed form of degree $(r-1,r-1)$
on $I$
\[
             W=\sum_{k=0}^{r-1}(z^*\xi)^{r-1-k}\nu^k.
\tag{I15.8}\label{r15i:W}
\]
It satisfies the literal identity of smooth forms
\[
                      (z^*\xi-\nu)W=(z^*\xi)^r-\nu^r.
\tag{I15.9}\label{r15i:telescoping}
\]

\begin{proposition}[The three actual incidence outputs]
\label{r15i:outputs}
Assume $(\mathrm{HP}_{\mathrm{nat}})_K$.
For every $\alpha\in\operatorname{PS}^{1,1}(E)$, put
$\beta_\alpha=W\wedge z^*\alpha\in
\operatorname{PS}^{r,r}(I)$.  The expressions
\[
 \mathcal G\alpha=p_*[g\beta_\alpha],\qquad
 \mathcal P\alpha=p_*[d'g\wedge\beta_\alpha],\qquad
 \mathcal Q\alpha=p_*[d''g\wedge\beta_\alpha]
\tag{I15.10}\label{r15i:three-outputs}
\]
have unique ordinary PWS representatives of degrees
$(0,0),(1,0),(0,1)$, with their prescribed pairings against
every ordinary smooth compact complementary test on $E$.
Here the displayed singular pairings are defined by the
following ordinary excision, not by an unproved product of
a current with a PWS form.
\end{proposition}

\begin{proof}
The two actual forms $z^*\alpha$ and $p^*\alpha$ restrict
equally to the algebraic closed immersion $\Sigma\hookrightarrow I$.
The frozen immersion-germ isomorphism, Round 14
\texttt{r14l:germ-iso}, makes them equal on an ordinary open
neighborhood $V$ of $\Sigma$.  Compactness permits a global
smooth cutoff $\chi$ that is zero near $\Sigma$ and one
outside $V$, with $\operatorname{Supp}(1-\chi)\subset V$.
Smooth ordinary partitions of unity supply this cutoff.

If $\eta$ is a complementary smooth test for respectively
the three outputs, the bidegrees of $\alpha\wedge\eta$ on
the $m$-dimensional target are respectively
$(m+1,m+1),(m,m+1),(m+1,m)$.
Thus on $V$
\[
             \beta_\alpha\wedge p^*\eta
                 =W\wedge p^*(\alpha\wedge\eta)=0.
\tag{I15.11}\label{r15i:cutoff-zero}
\]
This is an equality of actual forms before integration.
It continues to kill the paired expression after multiplication
by $g$, $d'g$, or $d''g$ off $\Sigma$.
Consequently all three prescribed pairings are unchanged by
inserting $\chi$ and are independent of its choice.

The forms $\chi g\beta_\alpha$,
$\chi d'g\wedge\beta_\alpha$ and
$\chi d''g\wedge\beta_\alpha$ extend by zero across a
whole ordinary neighborhood of $\Sigma$.  They are global
PWS forms, of degrees $(r,r),(r+1,r),(r,r+1)$.
Apply Theorem~\ref{r15ps:projective-bundle} successively to
$\rho$ (relative dimension one) and to $\lambda$ (relative
dimension $r-1$).  This constructs the claimed PWS outputs.
For a final smooth test on $E$, its pullback through each
proper stage remains smooth and compactly supported.  The
two prescribed pairing identities therefore compose without
ever using a PWS tensor as a test in a nonfinite pushforward.
The resulting outputs have exactly the uncut pairings above.
Current injectivity proves their uniqueness, cutoff independence,
and linearity in $\alpha$.
\end{proof}

\begin{theorem}[All-rank projective-bundle Aeppli formula]
\label{r15i:aeppli}
Assume $(\mathrm{HP}_{\mathrm{nat}})_K$, the frozen global
first-strip PWS replacement, proper scalar constancy, and the
frozen constant-fiber Aeppli pullback injection.  For every
smooth integral projective $S/K$ and every vector bundle $N$
of rank $r+1\geq2$,
\[
 H_{\mathrm A,\mathrm{act}}^{1,1}(\mathbf P_S(N))
      =\pi^*H_{\mathrm A,\mathrm{act}}^{1,1}(S)
                                      \oplus\mathbf R[\xi].
\tag{I15.12}\label{r15i:aeppli-formula}
\]
The class $[\xi]$ can equivalently be taken to be the actual
class of a model first Chern form for $\mathcal O_E(1)$.
For rank one the projective bundle is the base itself, with
the evident formula without an added summand.
\end{theorem}

\begin{proof}
Replace an arbitrary actual Aeppli class by a $D$-closed PWS
representative $\alpha\in\operatorname{PS}^{1,1}(E)$,
using Round 14's global first-strip replacement theorem.
The following outputs are defined by projective-bundle PWS
integration:
\[
 c_\alpha=\rho_*(z^*\alpha),\qquad
 J\alpha=\pi_*(\xi^r\wedge\alpha),\qquad
 \sigma=\lambda_*(\nu^r).
\tag{I15.13}\label{r15i:auxiliary-outputs}
\]
Their degrees are respectively $(0,0)$ on $F$ and $(1,1)$
on $S$ and $E$.  The first is $D$-closed and is a real
constant by proper scalar constancy on $F$.  The other two
are $D$-closed actual forms.  These facts follow by current
differentiation and current injectivity from $D\alpha=0$
and the separate closedness of the smooth Chern forms.
They do not feed derivatives of positive delta codimension
into the PWS pushforward theorem.

\emph{The all-smooth-test homotopy.}
Let $\eta\in\mathcal A_c^{m-1,m-1}(E)$ be arbitrary and
put $\Phi=\beta_\alpha\wedge p^*\eta$.
Since $\dim I=m+r$, this is an actual test of bidegree
$(m+r-1,m+r-1)$.  Moreover $D\beta_\alpha=0$.
The scalar actual Green identity \eqref{r15i:scalar-Green}
and \eqref{r15i:telescoping} give
\[
 \int_{I\setminus\Sigma}gD\Phi
  =\int_I\bigl((z^*\xi)^r-\nu^r\bigr)
                           \wedge z^*\alpha\wedge p^*\eta
        -\int_\Sigma W|_\Sigma\wedge\lambda^*(\alpha\eta).
\tag{I15.14}\label{r15i:scalar-pairing}
\]
The last integral is $\int_E\alpha\eta$.
Indeed
$W|_\Sigma=\sum_{k=0}^{r-1}(\lambda^*\xi)^{r-1-k}\nu^k$.
For $k<r-1$, multiplication by $\lambda^*(\alpha\eta)$
gives zero already as a form: its base bidegree exceeds $m$.
For $k=r-1$, apply the frozen projective-bundle degree
formula (R10.1), recalled in
Lemma~\ref{r15i:normalization}, to the compact actual top
test $\alpha\eta$ on $E$.  That formula explicitly permits
this nonsmooth test and gives
\[
        \int_F\nu^{r-1}\lambda^*(\alpha\eta)
                                          =\int_E\alpha\eta.
\]
Alternatively the same equality follows from the newly
proved PWS projection formula
\eqref{r15ps:PWS-projection} and compact PWS test comparison
of Corollary~\ref{r15ps:PWS-tests}, since
$\lambda_*(\nu^{r-1})=1$.  In neither route is a PWS test
inserted into a theorem that only permits smooth tests.

For explicit control of the boundary terms, on a neighborhood
of $\Sigma$ the forms
$\beta_\alpha p^*d'\eta$ and
$\beta_\alpha p^*d''\eta$ vanish by the same target-degree
calculation as \eqref{r15i:cutoff-zero}.
Their products with $g$ therefore extend by zero as global
actual forms.  Actual Stokes applies to them and gives
\[
 \begin{split}
 \big\langle D\mathcal G\alpha+d''\mathcal P\alpha
                         -d'\mathcal Q\alpha,\eta\big\rangle
 &=\int_I\bigl(g\beta_\alpha D\eta_p
       +d'g\,\beta_\alpha d''\eta_p
       -d''g\,\beta_\alpha d'\eta_p\bigr)\\
 &=-\int_{I\setminus\Sigma}gD(\beta_\alpha\eta_p),
 \qquad \eta_p=p^*\eta.
 \end{split}
\tag{I15.15}\label{r15i:boundary-sign}
\]
To check the signs independently of $r$, the total degree of
$\beta_\alpha$ is $2r$, hence
\[
 D(\beta_\alpha\eta_p)
   =-d''\beta_\alpha d'\eta_p
        +d'\beta_\alpha d''\eta_p+\beta_\alpha D\eta_p.
\]
Stokes for $g\beta_\alpha d''\eta_p$ and
$g\beta_\alpha d'\eta_p$ gives precisely
\eqref{r15i:boundary-sign}.  No differentiated cutoff term
is silently discarded.

\emph{Evaluation of the two curvature terms.}
For the first, the top actual PWS form
\[
       \Theta=\operatorname{pr}_z^*(\xi^r\alpha)
                              \wedge\operatorname{pr}_y^*\eta
\]
on $E\times_S E$ has zero restriction to the diagonal by
dimension.  The immersion-germ isomorphism makes it zero on
an ordinary neighborhood of that diagonal.  Its compact
support and that of $f^*\Theta$ are therefore on the
isomorphism locus of \eqref{r15i:incidence-map}.  Ordinary
change of variables on that locus gives
\[
 \int_I(z^*\xi)^r z^*\alpha\,p^*\eta
       =\int_{E\times_S E}\Theta
       =\int_E\pi^*(J\alpha)\wedge\eta.
\tag{I15.16}\label{r15i:first-term}
\]
The second equality is the projective-bundle base-change
identity for $\pi:E\to S$ pulled back along the algebraic
map $\pi:E\to S$, with its all-smooth-test pairing.
There is no actual pushforward along $f$ or $z$ in this step.

For the second curvature term use the tower $p=\lambda\rho$.
The target form $\nu^r\lambda^*\eta$ is smooth, so it is
an allowed test for the $\rho$-integration.  Thus
\[
 \int_I\nu^r z^*\alpha\,p^*\eta
       =\int_F c_\alpha\nu^r\lambda^*\eta
       =\int_E c_\alpha\sigma\wedge\eta.
\tag{I15.17}\label{r15i:second-term}
\]
The last equality uses only a smooth final test.  This records
why smooth metrics, rather than arbitrary model curvatures,
were chosen inside $W$.

Combining \eqref{r15i:scalar-pairing}--\eqref{r15i:second-term}
for every ordinary smooth compact test, and then using
actual-current injectivity, proves the actual identity
\[
 \alpha=\pi^*(J\alpha)-c_\alpha\sigma
          +D\mathcal G\alpha+d''\mathcal P\alpha
                                      -d'\mathcal Q\alpha.
\tag{I15.18}\label{r15i:homotopy}
\]

\emph{Eliminating the auxiliary class.}
Apply the same identity to $\alpha=\xi$.
The restriction of $z^*\mathcal O_E(1)$ to every
$\rho$-fiber is its hyperplane line.  Lemma~\ref{r15i:normalization}
gives $c_\xi=1$.
Subtract $c_\alpha$ times that identity from
\eqref{r15i:homotopy}.  It gives
\[
 \begin{split}
 \alpha-c_\alpha\xi
   &=\pi^*(J\alpha-c_\alpha J\xi)
       +D(\mathcal G\alpha-c_\alpha\mathcal G\xi)\\
   &\quad+d''(\mathcal P\alpha-c_\alpha\mathcal P\xi)
       -d'(\mathcal Q\alpha-c_\alpha\mathcal Q\xi).
 \end{split}
\tag{I15.19}\label{r15i:normalized-homotopy}
\]
The base form in parentheses is $D$-closed; the other three
terms are actual Aeppli boundaries.  This proves spanning in
\eqref{r15i:aeppli-formula} without identifying $\sigma$
with a Segre form or invoking any metric transgression for
a vector bundle.

For directness, restrict a relation
$\pi^*[\gamma]+a[\xi]=0$ to a projective line in a fiber
over a closed point of $S$ with residue field $F$.
The first term restricts to zero, and the second has integral
$a[F:K]$.  Actual Stokes kills the restricted Aeppli boundary,
so $a=0$.  Finally the frozen constant-fiber theorem,
Round 14 \texttt{r14j:equidimensional}, makes
$\pi^*:H_{\mathrm A}^{1,1}(S)\to H_{\mathrm A}^{1,1}(E)$
injective, proving $[\gamma]=0$.  All its hypotheses hold:
the map is projective dominant between normal integral
varieties with constant fiber dimension $r$, and the
Aeppli bidegree has minimum one.  The comparison with a
model Chern representative follows from the global PWS
metric-ratio scalar.
\end{proof}

\begin{corollary}[Degree-two Bott--Chern descent through every smooth center]
\label{r15i:all-smooth-centers}
Assume $(\mathrm{HP}_{\mathrm{nat}})_K$ and the frozen
foundations used above.  Let $X/K$ be smooth integral
projective, and let $C\hookrightarrow X$ be a smooth closed
nowhere dense center.  If $q:Y=\operatorname{Bl}_C X\to X$,
then
\[
          q^*:H_{\mathrm{BC,act}}^{2,2}(X)
                    \longrightarrow H_{\mathrm{BC,act}}^{2,2}(Y)
             \quad\hbox{is injective}.
\tag{I15.20}\label{r15i:all-center-BC22}
\]
There is no restriction on the dimension of $X$ or of $C$.
\end{corollary}

\begin{proof}
A codimension-one smooth center is Cartier and its blowup
is an isomorphism.  For an integral center of codimension
$c\geq2$, let $i:E_C\hookrightarrow Y$ be the exceptional
divisor and $\pi_C:E_C=\mathbf P_C(N_{C/X})\to C$.
The new normalized boundary quotient, proved separately in
Proposition~\ref{r15:degree-two-all-centers}, identifies the kernel with
\[
 \frac{H_{\mathrm A,\mathrm{act}}^{1,1}(E_C)}
 {i^*H_{\mathrm A,\mathrm{act}}^{1,1}(Y)
                  +\pi_C^*H_{\mathrm A,\mathrm{act}}^{1,1}(C)}.
\tag{I15.21}\label{r15i:exceptional-quotient}
\]
Theorem~\ref{r15i:aeppli} says that its numerator is the
sum of the base summand and the hyperplane class.
The hyperplane line is
$\mathcal O_Y(-E_C)|_{E_C}=\mathcal O_{E_C}(1)$.
Its actual Chern class therefore belongs to the first
denominator summand, and the quotient is zero.
If the smooth center has several connected components,
they are disjoint smooth integral projective centers.
Blowing them up successively gives the same blowup and
applies the integral-center conclusion component by
component.  This also covers different codimensions on
different components.  No assertion about finite behavior
of a general alteration on a prescribed cycle complement
enters this proof.
\end{proof}

\paragraph{Scope of this new consequence.}
The corollary controls degree-two Bott--Chern primitive descent
through all smooth-center modifications. This proof uses the
PWS projective-bundle output theorem. The later actual
projective-line theorem, Theorem~\ref{r15a2:projection},
enlarges the input class through smaller input degree two.
The full projective-space $(2,2)$ calculation consequently
adds point-center Bott--Chern $(3,3)$ descent, but higher
exceptional groups over positive-dimensional centers are
not controlled by the present argument.
Passing from the displayed primitive descent to new strong
Green theorems still requires a separately verified upstairs
curvature-obstruction vanishing theorem in the relevant cycle
degree.  In particular this section does not assert the
general strong Green conjecture or a general unconditional
Beilinson--Bloch height.
